%----------------------------------------------------------------------------------------
%	PACKETS AND CONFIGURATION
%----------------------------------------------------------------------------------------

\documentclass[12pt, a4paper]{article}

\usepackage{times} % Times New Roman font
\usepackage{graphicx} % 'graphics' package interface
\usepackage{geometry} % Edit document margins
\usepackage{hyperref} % Table of contents hyperlinks

% HELPER PACKAGES (REMOVE IN FINAL) %
\usepackage{blindtext} % Lorem Ipsum
\usepackage{todonotes} % TODOs as useful reminders
\setlength{\marginparwidth}{2cm} % Otherwise todonotes gets angry at me lol
\setuptodonotes{fancyline, color=green!40, shadow} % TODOs options
% HELPER PACKAGES (REMOVE IN FINAL) %

\graphicspath{ {./res/} } % Path to graphics
\hypersetup{    % ToC Hyperlink setup
    colorlinks,
    citecolor=black,
    filecolor=black,
    linkcolor=black,
    urlcolor=black
}

%----------------------------------------------------------------------------------------
%	DOCUMENT
%----------------------------------------------------------------------------------------

\begin{document}

\newgeometry{top=7cm, bottom=2cm} % Setting the margins for the title

% Title
\begin{titlepage}
    \centering
    {\Huge\bfseries Pandemic Information System Model\par} % Project title
    \vspace{1.5cm}
    {\scshape\large Systems and Methods for Big and Unstructured Data \par} % Course
    \vspace{0.5cm}
    {\scshape\large Prof. Marco Brambilla \par} % Professor
    \vspace{1cm}
    {\scshape\large % Description
        First delivery \par 
        Neo4j Project \par 
    }
    \vspace{0.5cm}
    {\slshape\large November 2021 \par} % Date
    \vspace{1cm}
    {\large\itshape % Authors
        \todo{add personal codes}
        Avci Oguzhan - \texttt{xxxxxxxx}\\
        Gentile Nicole - \texttt{10594355}\\
        Rigamonti Davide - \texttt{10629791}\\
        Singh Raul - \texttt{xxxxxxxx}\\
        Tagliaferri Mattia - \texttt{xxxxxxxx}
    }
    \vfill
    \begin{figure}[b]
        \includegraphics[scale=0.6]{polimi} % Polimi logo
        \centering
    \end{figure}

    \pagenumbering{gobble} % Remove page number

\end{titlepage}

\newgeometry{bottom=3cm} % Reset the margins
\pagenumbering{arabic} % Reset the page number

\clearpage

% INDEX
\tableofcontents

% LIST OF TODOs (REMOVE IN FINAL)
\listoftodos

\clearpage

% INTRODUCTION
\section{Introduction}

\subsection{Problem Specification}

We represented a population of XXX \todo{fill with population size} people in 
the USA. \\
The database receives data coming from tracing applications that use sensor in 
smartphones, wearable objects or other devices to understand whether two people 
had a contact; data includes date and time of the contact. \\  
Some places, like restaurants, theaters and hospitals, collect date and name of 
visiting people. \\
The idea is to use this database so that, if a person gets positive, we can 
understand who are all people who had a contact with him/her. Contacts are of 
different types: in family (since people from the same family are always in 
contact), in a location (if someone is positive, we alert people who were in 
the same location on the same day) or are given by an application using sensors. 
Data are recorded from 02/2020 and can be used for analytical purposes. 

\subsection{Hypoteses}

Vaccine date, date of the last contagion, date of the last negative test fields, 
healing date are optional fields: a person is positive if the contagion date is 
after the date of the last negative test; people can do tests without being 
infected; people can decide not to get the vaccine; vaccinated people can get 
infected; people who healed from covid can do tests after healing. \\
We also assumed that people can get covid at most once (realistic if we consider 
antibodies); in this way we can use data to build statistics. \\
Members of a family are all the people who live together, relatives who see 
each other very often or roommates. Obviously, all members of a family live in 
the same city. \\
When adding data about people going to places, we do not consider distances 
between the city where they live in and the location they go to. So, they can 
reach different places in different cities during the same day, assuming they 
do this in different moments. \\
Regarding places we have: people who got covid after 18/10 went to places from 
18/09 to 17/10, since we assume that later they are in quarantine; all the 
others went to visit locations from 18/09/2021 to 17/11/2021. We don’t retain 
previous data since it is useless for the queries. Moreover, if a person becomes 
positive, we alert all people who were in the same location in a span of time 
going from one hour earlier to two hours later than when the person went 
there. We assumed that hospitals also record people who are hospitalized for covid reasons and people can go to hospitals to make a visit (this is different 
from hospitalization). For simplicity, for people who are in hospital date of hospitalization and 
contagion coincide; we also assume they leave the hospital on the healing 
date. \\ If a person becomes positive, we alert all people who were in the same location in a span of time going from one hour earlier to two hours later than when the person went there.
In this way we alert: the ones who went 
there one hour earlier but may still be there and the ones who arrived after the 
infected and may have met him/her. \\
In order to populate the database, we imposed that every person visits from 6 
to 10 locations. 

\clearpage

% DATABASE
\section{Database}

\subsection{ER Diagram}

\blindtext

\subsection{Dataset description}

We used three types of nodes: Person, Location and City. 
People are characterized by their name and surname, birthdate, city, email, social security number, vaccine date, date of the last contagion, date of the last negative test and healing date. People can be linked by family relationships.
Locations have a name and a type (restaurant, theater, hospital); type is used so that the dataset can be easily expanded with new types of location.
Possible relationships are: 
\begin{itemize}
	\item \texttt{WENT\char`_TO} to track people who visited a location at a certain time;
	\item \texttt{IN\char`_FAMILY} for family relationships;
	\item \texttt{IS\char`_IN} between locations and cities;
	\item \texttt{LIVES\char`_IN} to link people to the city they live in;
	\item \texttt{HAS\char`_MET} to indicate contacts between people using tracing app or devices;
	\item \texttt{IS\char`_HOSPITALIZED\char`_IN} to indicate people who are/were hospitalized for covid reasons.
\end{itemize}

\subsection{Queries}

\blindtext

\subsection{Commands}

\blindtext

\clearpage

% APPLICATION
\section{Application}

\subsection{Description}

\blindtext

\subsection{User Guide}

\blindtext

\clearpage

% REFERENCES AND SOURCES
\section{References and sources}

\blindtext

\clearpage

\end{document}